% Shared visual style for diagrams and algorithms.
% Included once from main.tex; not a chapter, contains no \chapter command.

\usepackage{tikz}
\usetikzlibrary{arrows.meta,positioning,calc,decorations.pathreplacing,patterns,matrix,fit,backgrounds,shapes.geometric}
\usepackage{pgfplots}
\pgfplotsset{compat=1.18}
\usepackage{float}
\usepackage{subcaption}
\usepackage{algorithm}
\usepackage{algpseudocode}

% ---- Palette: restrained, print-friendly, colorblind-conscious ----
\definecolor{pqcblue}{HTML}{1F4E79}   % primary structure, arrows, key vectors
\definecolor{pqcteal}{HTML}{2E7D6B}   % secondary structure
\definecolor{pqcorange}{HTML}{C1601B} % emphasis / the "answer" or highlighted object
\definecolor{pqcgray}{HTML}{6B6B6B}   % grid lines, muted context
\definecolor{pqclight}{HTML}{EDEDED}  % fills / backgrounds

% ---- Reusable TikZ styles ----
\tikzset{
  latticept/.style={circle, fill=pqcgray, inner sep=0pt, minimum size=3pt},
  latticeptbold/.style={circle, fill=pqcblue, inner sep=0pt, minimum size=4pt},
  basisvec/.style={-{Latex[length=2.5mm]}, pqcblue, thick},
  altvec/.style={-{Latex[length=2.5mm]}, pqcteal, thick},
  shortvec/.style={-{Latex[length=2.5mm]}, pqcorange, thick, line width=1pt},
  flownode/.style={draw=pqcblue, thick, rounded corners=2pt, fill=pqclight,
                    minimum height=7mm, inner sep=3pt, align=center, font=\small},
  flowarrow/.style={-{Latex[length=2.5mm]}, pqcgray, thick},
  keynode/.style={draw=pqcorange, thick, rounded corners=2pt, fill=white,
                   minimum height=7mm, inner sep=3pt, align=center, font=\small},
}

% ---- Algorithm captions read "Algorithm N: Name" consistently ----
\floatname{algorithm}{Algorithm}
\renewcommand{\algorithmicrequire}{\textbf{Input:}}
\renewcommand{\algorithmicensure}{\textbf{Output:}}

% ---- Theorem-like environments (shared counter, numbered by chapter) ----
\usepackage{amsthm}
\theoremstyle{definition}
\newtheorem{definition}{Definition}[chapter]
\newtheorem{example}[definition]{Example}
\theoremstyle{plain}
\newtheorem{theorem}[definition]{Theorem}
\newtheorem{proposition}[definition]{Proposition}
\newtheorem{lemma}[definition]{Lemma}
\theoremstyle{remark}
\newtheorem*{remark}{Remark}

% ---- Running headers show "Chapter N. Title" on every page, not section
% titles. Long \section{} titles otherwise collide with the page number
% when they land on a recto page; showing chapter context on both sides
% is standard for a book this size and sidesteps that failure mode
% entirely. Individual chapters can still shorten their own header with
% \chaptermark{...} right after \chapter{...} (see e.g. 14_cbd.tex).
\makeatletter
\renewcommand{\sectionmark}[1]{}
\renewcommand{\chaptermark}[1]{%
  \markboth{\MakeUppercase{\ifnum \c@secnumdepth >\m@ne \@chapapp\ \thechapter.\ \fi #1}}%
           {\MakeUppercase{\ifnum \c@secnumdepth >\m@ne \@chapapp\ \thechapter.\ \fi #1}}%
}
\makeatother

% Reference to a companion Sage file, hyperlinked to the repository.
\newcommand{\sagerepo}{https://github.com/AppliedPQC/AppliedPQC/blob/main/sage}
\newcommand{\sagefile}[1]{\href{\sagerepo/#1}{\texttt{sage/#1}}}

% Pointer to the browser-runnable version of the companion Sage code.
\newcommand{\playgroundurl}{https://appliedpqc.io/playground.html}
\newcommand{\playground}[1]{%
  \smallskip\noindent\emph{Run it in a browser:} #1. See
  \href{\playgroundurl}{appliedpqc.io/playground.html} --- nothing to install.}

